\documentclass[../cours_1re.tex]{subfiles}
\begin{document}
\setcounter{chapter}{10}
\chapter{Optimiser, comparer et approcher}
\label{ch:11}
\begin{objectif}
Résoudre une optimisation avec ses contraintes, établir une inégalité et comprendre Newton. Prérequis : dérivation et tableaux de variations.
\end{objectif}
\section{Du problème à la fonction}
Optimiser exige trois éléments : une quantité à maximiser ou minimiser, un domaine admissible, et une justification de l'optimum. Une valeur fournie par un logiciel est un candidat.
\begin{exemple}[Une boîte sans couvercle]
On découpe aux quatre coins d'une feuille carrée de côté $6$ des carrés de côté $x$, puis on replie. Le volume est
$V(x)=x(6-2x)^2=4x(3-x)^2$, avec $0<x<3$.
La dérivée est $V'(x)=12(x-1)(x-3)$. Elle est positive sur $]0;1[$, négative sur $]1;3[$.
Le maximum est atteint en $x=1$ et vaut $16$ unités de volume. Les limites nulles aux bornes correspondent à des boîtes dégénérées, exclues du domaine mais utiles pour interpréter le résultat.
\end{exemple}
\section{Comparer deux courbes}
Étudier la position de $f$ par rapport à $g$ revient à étudier le signe de $h=f-g$. Les zéros sont les abscisses d'intersection ; $h>0$ signifie que le graphe de $f$ est au-dessus.
\begin{exemple}
Pour $x>0$, montrer $x+1/x\ge2$. Poser $h(x)=x+1/x-2$.
Alors $h'(x)=1-1/x^2$, négative avant $1$ et positive après ; le minimum est $h(1)=0$.
Une seconde preuve est l'identité $h(x)=(x-1)^2/x\ge0$. Choisir l'algèbre quand elle est plus courte.
\end{exemple}
\section{Une tangente comme borne}
\approfondissement Si $f'$ est croissante sur un intervalle, le graphe est au-dessus de chacune de ses tangentes :
\[
f(x)\ge f(a)+f'(a)(x-a).
\]
\begin{demonstration}
Poser $g(x)=f(x)-f(a)-f'(a)(x-a)$. Alors $g'(x)=f'(x)-f'(a)$ est négative ou nulle avant $a$, positive ou nulle après. $g$ possède donc son minimum en $a$, où $g(a)=0$.
\end{demonstration}
Cette propriété est une caractérisation de la convexité pour les fonctions dérivables. Pour $f(x)=x^2$, l'écart à la tangente en $a$ est exactement $(x-a)^2$. Pour une dérivée décroissante, les inégalités s'inversent : c'est la concavité.
\section{Complément : une erreur d'ordre deux}
\begin{theoreme}
Si $f$ possède une dérivée seconde sur un voisinage du segment joignant $a$ à $a+h$ et si $|f''|\le M$ sur ce segment, alors
\[
|f(a+h)-f(a)-hf'(a)|\le \frac M2 h^2.
\]
\end{theoreme}
\begin{demonstration}
Pour $h\ne0$, poser $R=f(a+h)-f(a)-hf'(a)$ et
$g(t)=f(t)-f(a)-(t-a)f'(a)-R(t-a)^2/h^2$.
On a $g(a)=g(a+h)=0$ et $g'(a)=0$. Deux applications successives de Rolle donnent un point $\xi$ strictement entre les bornes où $g''(\xi)=0$. Ainsi $R=f''(\xi)h^2/2$. Pour $h=0$, l'affirmation est immédiate.
\end{demonstration}
Cette preuve utilise Rolle, dont les fondements sont explicités aux chapitres 10 et 25. Une petite distance à $a$ ne suffit pas seule : la borne sur $f''$ contrôle la courbure.
\section{Newton : annuler la tangente}
La tangente à $f$ en $u_n$ coupe l'axe horizontal en
\[
u_{n+1}=u_n-\frac{f(u_n)}{f'(u_n)},\qquad f'(u_n)\ne0.
\]
C'est la méthode de Newton : on remplace localement l'équation par une équation affine. Une division possible ne garantit pas que l'itération converge.
\begin{exemple}[Un cas favorable : racine carrée]
Pour chercher $\sqrt A$, $A>0$, prendre $f(x)=x^2-A$ et $u_0>0$ :
$u_{n+1}=(u_n+A/u_n)/2$.
Pour $A=2$, $u_0=2$, on obtient $u_1=3/2$, $u_2=17/12$, $u_3=577/408$.
\end{exemple}
\begin{demonstration}[Convergence et erreur]
Noter $r=\sqrt A$. Tout terme reste positif et
\[
u_{n+1}-r=\frac{(u_n-r)^2}{2u_n}\ge0.
\]
Donc, dès le premier terme calculé, $u_n\ge r$. Alors
$u_{n+1}-u_n=(A-u_n^2)/(2u_n)\le0$ : la suite décroît et est minorée. Elle converge vers un réel $\ell\ge r>0$ ; passer à la limite donne $\ell=(\ell+A/\ell)/2$, donc $\ell=r$.
De plus, pour $n\ge1$, $0\le u_{n+1}-r\le(u_n-r)^2/(2r)$ : l'erreur est quadratique.
\end{demonstration}
On obtient aussi un encadrement certifié $A/u_n\le r\le u_n$ pour $n\ge1$. Sa largeur fournit un critère d'arrêt vérifiable.
\begin{attention}
Pour $f(x)=x^3-2x+2$, Newton démarrant en $0$ passe à $1$, puis revient à $0$. Les dérivées utilisées ne sont pas nulles, pourtant la méthode boucle. Le choix du départ et les hypothèses de convergence sont essentiels.
\end{attention}
\section{Une grille ne prouve pas un optimum}
Un balayage teste seulement les points d'une grille. Pour $f(x)=1-(x-0{,}3)^2$ sur $[0;1]$, une grille $0,0{,}5,1$ trouve $0{,}96$, alors que le vrai maximum vaut $1$ en $0{,}3$.
Si l'on sait que $|f'|\le L$ et si tout point du domaine est à distance au plus $\delta$ d'un point testé, les accroissements finis donnent une erreur sur le maximum au plus $L\delta$. La garantie vient de cette estimation, pas de la seule densité visuelle du graphe.

\begin{finchapitre}
\section{Utilité et prolongements}
Optimisation, inégalités et recherche de racines partagent les mêmes outils : domaine, dérivée, signe et contrôle de l'erreur.

En licence puis en master, la convexité permet de garantir qu'un minimum local est global. Newton devient une méthode de résolution de systèmes grâce à la matrice jacobienne. La distinction entre convergence observée et convergence prouvée reste déterminante en calcul scientifique.
\begin{retenir}
Un résultat d'optimisation doit préciser où il est atteint. Une méthode numérique doit préciser ses hypothèses, son critère d'arrêt et ce que garantit effectivement l'erreur annoncée.
\end{retenir}
\end{finchapitre}
\end{document}
